\documentclass[UTF8,a4paper,11pt]{ctexart}

\usepackage[margin=2.4cm]{geometry}
\usepackage{amsmath,amssymb}
\usepackage{booktabs}
\usepackage{array}
\usepackage{enumitem}
\usepackage{hyperref}
\usepackage{xcolor}
\usepackage{longtable}

\hypersetup{
  colorlinks=true,
  linkcolor=blue!60!black,
  citecolor=blue!60!black,
  urlcolor=blue!60!black
}

\newcommand{\todo}[1]{\textcolor{red!70!black}{[待补充：#1]}}
\newcommand{\tok}[1]{\ensuremath{\mathrm{#1}}}
\newcommand{\yes}{\ensuremath{\checkmark}}

\title{FastWAM + CoT DiT：基于思维链的世界--动作模型实现计划初稿}
\author{上海创智博士夏令营三人组队项目}
\date{\today}

\begin{document}
\maketitle

\begin{abstract}
本文档给出在 FastWAM 代码库上实现 WAM-CoT 的初步技术计划。核心思想是在已有 Video DiT 与 Action DiT 的 MoT 框架中新增一个 CoT DiT 分支：其输入为由 Qwen3-VL-2B 提取的 VLM tok，结构参考 Motus 的 Understanding Expert，并通过统一的 MoT attention mask 与 Video DiT、Action DiT 交互。该设计希望在保留 FastWAM 高效动作预测能力的同时，引入来自 VLM 的场景理解、空间关系与任务语义先验，为长视野、遮挡和多阶段操作任务提供内部物理思维链表征。
\end{abstract}

\section{项目目标与背景}

课题目标是研究基于思维链机制的世界--动作模型（WAM-CoT）是否能够提升机器臂在复杂物理交互任务中的策略性能，尤其关注遮挡、长视野、多阶段堆叠等容易出现误差累积和策略崩溃的任务。本文选择“内部物理思维链”路线：不把 CoT 仅作为外部文本规划，而是在 WAM 内部加入可与视频 token 和动作 token 共同推理的 latent reasoning token。

FastWAM 当前以 Wan2.2 为 Video DiT backbone，额外构造 Action DiT 分支，并通过 Mixture-of-Transformers（MoT）让 action token 与输入帧的 video VAE latent 交互。训练时 Video DiT 预测未来视频噪声，Action DiT 预测未来动作噪声；推理时可跳过未来视频生成，仅基于首帧和语言/状态条件预测动作，从而提高策略执行效率。

Motus 同样基于 Wan2.2，但进一步引入 Qwen3-VL-2B 和 Understanding Expert，使 VLM 的语义理解 token 参与 Tri-model Joint Attention。本文计划将这一思路迁移到 FastWAM，形成包含 Video DiT、Action DiT、CoT DiT 三个专家的 FastWAM+CoT DiT。

\section{已有基础}

\subsection{FastWAM 的两专家 MoT}

FastWAM 的核心专家包括：
\begin{itemize}[leftmargin=1.6em]
  \item Video DiT：继承 Wan2.2 TI2V-5B 的视频生成结构。输入为 VAE latent，经 patchify 后形成 video token；其 hidden size 为 3072，30 层，24 个 attention head，每个 head 128 维。
  \item Action DiT：与 Video DiT 层数、head 数、head 维度保持一致，hidden size 为 1024，FFN 维度为 4096。输入为动作序列，经 action encoder 映射为 action token，并用独立 flow timestep 做 AdaLN 调制。
  \item MoT：每一层分别由各 expert 产生自己的 Q/K/V，然后按 expert 顺序拼接，在同一个 attention mask 下做 mixed self-attention，最后再按 expert 切分输出并接回各自的 cross-attention、FFN 和 head。
\end{itemize}

当前 FastWAM 的基础 mask 逻辑为：Video DiT 内部使用 \texttt{first\_frame\_causal}，即首帧 token 只能看首帧，未来 video token 可以看首帧和未来 video token；Action DiT 内部 token 之间全连接；Action DiT 只能看首帧 video token。这对应“动作只依赖当前观测，不偷看未来视频”的策略设定。

\subsection{Motus 的三专家思想}

公开 Motus 资料显示，其 MoT 由三类专家组成：Video Generation Model（Wan2.2-5B）、Action Expert 和 Understanding Expert。三者通过 Tri-model Joint Attention 共享多头自注意力的信息流，但各自保留独立的 normalization、FFN、输入/输出适配器和 timestep 处理逻辑。Motus 的设计动机是把视频生成的物理先验、动作建模的控制先验、VLM 的语义理解先验放在同一个联合生成框架内。

Motus 中 Action Expert 的公开实现要点如下：
\begin{itemize}[leftmargin=1.6em]
  \item 输入编码：finetune 模式下分别编码 state token 和 action token，再拼接；pretrain 模式下可只编码 action token。编码器使用 \texttt{mlp3x\_silu}，并加入固定一维正余弦位置编码。
  \item 主干结构：层数与 Wan 一致。每个 block 包含 LayerNorm、AdaLN timestep 调制、FFN，以及映射到 Wan attention head space 的 action Q/K/V 投影。
  \item MoT 交互：Action Expert 不直接复用 Wan 的 hidden size，而是用参数 \texttt{wan\_action\_qkv} 将 action hidden 映射到 Wan 的 \texttt{num\_heads $\times$ head\_dim} 空间，参与 joint attention 后再通过 \texttt{wan\_action\_o} 投回 action hidden size。
  \item 输出头：Action Decoder 使用 timestep 调制后的 LayerNorm 和 MLP 输出动作噪声/速度场。
\end{itemize}

Motus 中 Understanding Expert 的公开实现要点如下：
\begin{itemize}[leftmargin=1.6em]
  \item 输入来自 VLM 最后一层 hidden states。公开配置中 VLM 为 Qwen3-VL-2B-Instruct，VLM hidden size 为 2048。
  \item VLM 输出先经过 \texttt{mlp3x\_silu} adapter，从 2048 维投影到 Understanding Expert hidden size，公开 RobotWin 配置中为 512。
  \item Understanding Expert 与 Action Expert 类似，也包含与 Wan 层数对齐的 Transformer blocks，并用 \texttt{wan\_und\_qkv} 映射到 Wan head space 参加 Tri-model Joint Attention。
  \item 与 Action Expert 不同，Understanding Expert 通常没有动作输出 decoder，也不需要 registers；其作用是更新并传递 VLM 语义 token，让 video/action 分支在每层 attention 中读取语义理解表征。
\end{itemize}

\section{VLM Tok 的获取方式}

为严格对齐 Motus，建议将 VLM tok 的产生过程写成独立模块，例如 \texttt{VLMTokenExtractor}，并在数据预处理或训练 forward 中复用同一逻辑。具体步骤如下：

\begin{enumerate}[leftmargin=1.8em]
  \item 构造 Qwen3-VL 对话输入。输入包括当前 RGB 图像和任务语言指令，格式为一条 user message：图像在前，文本指令在后。
  \item 调用 Qwen processor 的 \texttt{apply\_chat\_template(..., add\_generation\_prompt=True)}，再通过 \texttt{process\_vision\_info} 处理图像，最终得到 \texttt{input\_ids}、\texttt{attention\_mask}、\texttt{pixel\_values}、\texttt{image\_grid\_thw}。
  \item 使用 Qwen3-VL 的 embedding 层生成 \texttt{inputs\_embeds}，调用 VLM 的 image encoder 得到 image features，并用 placeholder mask 将 image features 写回对应 image token 位置。
  \item 根据 \texttt{input\_ids}、\texttt{image\_grid\_thw} 和 \texttt{attention\_mask} 调用 Qwen3-VL 的 RoPE index 逻辑生成 \texttt{position\_ids}。
  \item 冻结 VLM 参数，用 \texttt{output\_hidden\_states=True} 前向 Qwen3-VL language model，取 \texttt{hidden\_states[-1]} 作为 VLM last-layer token。
  \item 通过 CoT/Understanding adapter 将 VLM hidden states 投影到 CoT DiT hidden size，得到本文定义的 VLM tok。
\end{enumerate}

形式化地，设 Qwen3-VL 最后一层 hidden states 为
\[
H^{\mathrm{vlm}} = \mathrm{Qwen3VL}_{\mathrm{frozen}}(I_0, \ell)_{\mathrm{last}}
\in \mathbb{R}^{B \times L_u \times d_{\mathrm{vlm}}},
\]
其中 $I_0$ 是当前观测图像，$\ell$ 是语言指令。VLM tok 为
\[
U_0 = \phi_{\mathrm{adapter}}(H^{\mathrm{vlm}})
\in \mathbb{R}^{B \times L_u \times d_u},
\]
其中 $\phi_{\mathrm{adapter}}$ 建议使用 Motus 一致的 \texttt{mlp3x\_silu}。

如果本项目内部称“Qwen-3-2B”，实现上建议以 Motus 配置中的 \texttt{Qwen3-VL-2B-Instruct} 为准；若实际使用纯文本 Qwen3-2B，则无法提供图像 token，和 Motus 的 VLM tok 定义不完全一致，需要另行改写输入语义。

\section{CoT DiT 的网络架构设计}

\subsection{设计原则}

CoT DiT 应满足三个约束：
\begin{itemize}[leftmargin=1.6em]
  \item 与 FastWAM 的 MoT 接口兼容：必须提供 \texttt{tokens}、\texttt{freqs}、\texttt{t\_mod} 或等价字段，并包含与 Video/Action 对齐的 \texttt{blocks}、\texttt{num\_heads}、\texttt{attn\_head\_dim}。
  \item 与 Motus Understanding Expert 对齐：输入来自 Qwen3-VL last hidden states，经 adapter 后参与每层 joint attention；不直接预测动作。
  \item 与 FastWAM 的计算预算匹配：CoT DiT 不应显著破坏 FastWAM 的推理效率，因此 hidden size 和 token 数要可控。
\end{itemize}

\subsection{推荐配置}

本文建议第一版使用如下配置：

\begin{center}
\begin{tabular}{lll}
\toprule
模块 & 推荐值 & 说明 \\
\midrule
VLM backbone & Qwen3-VL-2B-Instruct & 与 Motus 对齐，训练时冻结 \\
VLM input dim & 2048 & Qwen3-VL-2B last hidden size \\
CoT hidden size $d_u$ & 512 & 参考 Motus Understanding Expert，节省显存 \\
CoT FFN dim & 2048 & 4 倍 hidden size \\
层数 & 30 & 与 Wan2.2 / Action DiT 对齐 \\
attention heads & 24 & 与 MoT mixed attention 对齐 \\
head dim & 128 & 与 Video/Action attention head space 对齐 \\
adapter & \texttt{mlp3x\_silu} & 2048 $\rightarrow$ 512 \\
输出 head & 无或 lightweight probe & 默认不直接监督输出 \\
\bottomrule
\end{tabular}
\end{center}

CoT DiT block 可采用：
\[
\begin{aligned}
\tilde{U}^{(l)} &= \mathrm{LN}_1(U^{(l)}),\\
Q_u,K_u,V_u &= \mathrm{Proj}^{u\rightarrow \mathrm{WanHead}}(\tilde{U}^{(l)}),\\
Y_v,Y_a,Y_u &= \mathrm{MoTAttn}\bigl([Q_v,Q_a,Q_u], [K_v,K_a,K_u], [V_v,V_a,V_u]; M\bigr),\\
U^{(l+\frac12)} &= U^{(l)} + \mathrm{Proj}^{\mathrm{WanHead}\rightarrow u}(Y_u),\\
U^{(l+1)} &= U^{(l+\frac12)} + \mathrm{FFN}(\mathrm{LN}_2(U^{(l+\frac12)})).
\end{aligned}
\]

与 Action DiT 的区别是：CoT DiT 不需要 action noise scheduler，也不需要输出动作维度；如果希望让 CoT token 也具备 denoising 式 latent reasoning，可加入一个固定或可学习的 CoT timestep embedding，但第一版建议不引入额外 diffusion loss，先把 CoT DiT 作为语义--物理中间表征分支。

\section{Attention Mask 设计与实现}

\subsection{符号}

本文把 token 分组为：
\begin{itemize}[leftmargin=1.6em]
  \item $u$：VLM tok / CoT token；
  \item $f_0$：首帧 video latent token；
  \item $f_1, f_h$：未来 video latent token，其中 $f_1$ 可表示第一组未来 token，$f_h$ 表示后续 horizon token；
  \item $a_1, a_h$：未来动作 token。
\end{itemize}

在实现中无需按 $f_1/f_h$ 单独存储，只需要知道每帧 video token 数 \texttt{video\_tokens\_per\_frame}，即可切出 $f_0$ 和未来 video token。

\subsection{训练 mask}

训练时的目标 mask 为：

\begin{center}
\begin{tabular}{c|cccccc}
\toprule
Query $\backslash$ Key & VLM tok & $f_0$ & $f_1$ & $f_h$ & $a_1$ & $a_h$ \\
\midrule
VLM tok & \yes & \yes & & & & \\
$f_0$   & \yes & \yes & & & & \\
$f_1$   & \yes & \yes & \yes & \yes & & \\
$f_h$   & \yes & \yes & \yes & \yes & & \\
$a_1$   & \yes & \yes & & & \yes & \yes \\
$a_h$   & \yes & \yes & & & \yes & \yes \\
\bottomrule
\end{tabular}
\end{center}

该 mask 的含义是：
\begin{itemize}[leftmargin=1.6em]
  \item VLM tok 只读取自身和当前首帧，避免被未来视频或动作标签污染；
  \item Video token 可读取 VLM tok，从语义理解中获得物体、空间关系和任务目标信息；
  \item 未来 video token 仍维持 FastWAM 原有 video-to-video 可见性；
  \item Action token 可读取 VLM tok 和首帧 token，但不可读取未来 video token，避免训练时从未来视觉答案中泄漏动作信息；
  \item Action token 内部全连接，维持 action chunk 的整体预测。
\end{itemize}

\subsection{推理 mask}

推理只保留当前观测和 action token，因此目标 mask 为：

\begin{center}
\begin{tabular}{c|cccc}
\toprule
Query $\backslash$ Key & VLM tok & $f_0$ & $a_1$ & $a_h$ \\
\midrule
VLM tok & \yes & \yes & & \\
$f_0$   & \yes & \yes & & \\
$a_1$   & \yes & \yes & \yes & \yes \\
$a_h$   & \yes & \yes & \yes & \yes \\
\bottomrule
\end{tabular}
\end{center}

这与 FastWAM 当前高效动作推理保持一致：不需要在测试时生成未来视频，只需 prefill $u$ 和 $f_0$，再 denoise action token。

\subsection{代码实现草案}

FastWAM 当前的 \texttt{\_build\_mot\_attention\_mask(video\_seq\_len, action\_seq\_len, video\_tokens\_per\_frame, device)} 返回 $[S,S]$ bool mask。加入 CoT DiT 后建议改为：

\begin{verbatim}
def _build_cot_mot_attention_mask(
    cot_seq_len: int,
    video_seq_len: int,
    action_seq_len: int,
    video_tokens_per_frame: int,
    device: torch.device,
) -> torch.Tensor:
    total = cot_seq_len + video_seq_len + action_seq_len
    mask = torch.zeros((total, total), dtype=torch.bool, device=device)

    u0, u1 = 0, cot_seq_len
    v0, v1 = u1, u1 + video_seq_len
    a0, a1 = v1, total
    first = min(video_tokens_per_frame, video_seq_len)

    # cot -> cot + f0
    mask[u0:u1, u0:u1] = True
    mask[u0:u1, v0:v0 + first] = True

    # video -> cot + original video mask
    mask[v0:v1, u0:u1] = True
    mask[v0:v1, v0:v1] = self.video_expert.build_video_to_video_mask(...)

    # action -> cot + f0 + action
    mask[a0:a1, u0:u1] = True
    mask[a0:a1, v0:v0 + first] = True
    mask[a0:a1, a0:a1] = True
    return mask
\end{verbatim}

MoT 的 expert 顺序应固定为 \texttt{\{"cot": cot\_expert, "video": video\_expert, "action": action\_expert\}}，从而 mask 的行列顺序与 token 拼接顺序一致。若为了最小改动保留 \texttt{video, action} 顺序，也可以将 CoT 追加到末尾，但 mask 构造必须同步调整。

\section{FastWAM+CoT DiT 完整 Pipeline}

\subsection{训练流程}

\begin{enumerate}[leftmargin=1.8em]
  \item 从样本中读取多相机/单相机视频、动作序列、语言指令和 proprioception。
  \item 使用 Video VAE 编码视频，得到 latent video；首帧作为条件，未来帧加入 flow-matching 噪声。
  \item 对动作序列加入 action scheduler 噪声，得到 noisy action token。
  \item 使用 Qwen3-VL-2B 提取当前图像和语言指令的 last hidden states，经过 adapter 得到 VLM tok。
  \item 分别执行 \texttt{video\_expert.pre\_dit}、\texttt{action\_expert.pre\_dit}、\texttt{cot\_expert.pre\_dit}，得到三类 token、RoPE 频率和调制参数。
  \item 构建训练 attention mask，将 VLM tok、video token、action token 拼接后送入三专家 MoT。
  \item Video DiT head 输出 video velocity/noise prediction；Action DiT head 输出 action velocity/noise prediction；CoT DiT 默认不设主损失。
  \item 总损失第一版保持
  \[
  \mathcal{L}=\lambda_v\mathcal{L}_{video}+\lambda_a\mathcal{L}_{action}.
  \]
  后续可加入辅助损失，如 VLM token dropout consistency、阶段成功预测 probe 或 object relation probe。
\end{enumerate}

\subsection{推理流程}

\begin{enumerate}[leftmargin=1.8em]
  \item 输入当前观测图像、语言指令和 proprioception。
  \item VAE 编码首帧得到 $f_0$；Qwen3-VL 提取 VLM tok 得到 $u$。
  \item 初始化 action latent 为高斯噪声。
  \item 每个 denoising step 中，构建推理 mask，让 action token 读取 $u$、$f_0$ 和 action token 自身。
  \item Action DiT 输出动作噪声/速度场，scheduler 更新 action latent。
  \item 完成所有 step 后反归一化动作 chunk，并交给仿真环境执行。
\end{enumerate}

\section{实验与消融计划占位}

后续实验建议至少包含：
\begin{itemize}[leftmargin=1.6em]
  \item Baseline：原始 FastWAM；
  \item Ours：FastWAM+CoT DiT；
  \item 去除 CoT：保留参数量相近的空 expert 或关闭 VLM tok；
  \item 打乱 VLM 输入：打乱语言或替换图像，验证 CoT token 是否真实参与动作生成；
  \item mask 消融：允许 action 看未来 video token、禁止 action 看 VLM tok、禁止 video 看 VLM tok；
  \item 任务维度：遮挡抓取、多阶段堆叠、长视野搬运等。
\end{itemize}

主要指标包括任务成功率、阶段成功率、平均执行步数、失败类型统计、推理时延、显存占用和 action chunk 稳定性。\todo{补充具体任务、数据规模、训练超参和评估脚本路径。}

\section{实现里程碑}

\begin{enumerate}[leftmargin=1.8em]
  \item M1：复现 FastWAM baseline 训练/推理，确认 RoboTwin 或 LIBERO 指标可跑通。
  \item M2：加入 Qwen3-VL VLM tok 提取脚本，支持离线缓存，降低训练时显存和吞吐压力。
  \item M3：实现 CoT DiT / Understanding Expert，并接入三专家 MoT。
  \item M4：实现训练和推理 attention mask，完成 shape test 和 overfit 小样本测试。
  \item M5：完整训练 FastWAM+CoT DiT，进行消融实验和失败案例分析。
\end{enumerate}

\section{当前风险与待确认事项}

\begin{itemize}[leftmargin=1.6em]
  \item 本机当前不可访问用户提供的私有 Motus 路径 \texttt{/pfs/pfs-iQ14no/yuhao/Motus}，本文 Motus 细节基于公开论文、README 和公开仓库代码整理。若私有版本有改动，需要以后按源码校正。
  \item CoT DiT 引入 Qwen3-VL，会显著增加显存占用。第一版应优先支持 VLM tok 离线缓存，训练时冻结 VLM。
  \item VLM token 长度可能较大，需要考虑截断、只保留 image/text 有效 token、或使用 pooling/query compression。
  \item 如果 Action token 读取 VLM tok 后出现训练不稳定，可先冻结 CoT DiT 前若干 epoch，仅训练 adapter 和 Action/Video 侧。
\end{itemize}

\begin{thebibliography}{9}
\bibitem{fastwam}
Fast-WAM: Do World Action Models Need Test-time Future Imagination? arXiv:2603.16666.

\bibitem{motus}
Hongzhe Bi et al. Motus: A Unified Latent Action World Model. arXiv:2512.13030.

\bibitem{wan}
Wan-AI. Wan2.2-TI2V-5B model family.

\bibitem{qwen3vl}
Qwen Team. Qwen3-VL-2B-Instruct.

\bibitem{cotvla}
CoT-VLA: Visual Chain-of-Thought Reasoning for Vision-Language-Action Models. arXiv:2503.22020.
\end{thebibliography}

\end{document}
